\chapter{Resultados}
\label{ch:resultados}

Este capítulo reporta las salidas de cada etapa para la Zona Metropolitana de
Guadalajara (\zmg). El orden de las secciones sigue la metodología: la superficie de
demanda estimada (\cref{sec:res-w1}), su calibración (\cref{sec:res-w2}), el
diagnóstico de brecha de cobertura y su interpretación supervisada
(\cref{sec:res-w3}), la priorización Nodo--Lugar--Personas (\cref{sec:res-w4}), los
corredores generados (\cref{sec:res-w6}), la auditoría de la red existente
(\cref{sec:res-w7}) y el conjunto de validación (\cref{sec:res-w8}).

\section{Superficie de demanda de transporte estimada (W1)}
\label{sec:res-w1}

La metodología de Nivel~1 produce una estimación de demanda de transporte para los
\num{1881} \ageb{}s, con un total de \num{8471576} extremos de viaje en transporte al
día. La propiedad vehicular de los hogares varía ampliamente en la metrópoli: la tasa
vehicular media es \num{0.577}, lo que da una propensión media al transporte de
\num{0.423}. Como la demanda se pondera por $(1-v_i)$ (\cref{eq:transit-demand}), los
\ageb{}s con alta propiedad automovilística aportan proporcionalmente menos demanda
latente aun donde la generación bruta de viajes es alta. La \cref{fig:demand} mapea la
superficie resultante, que se concentra a lo largo de los densos corredores central y
oriente de la metrópoli.

\begin{figure}[htbp]
  \centering
  \includegraphics[width=0.82\linewidth]{zmg_demand.pdf}
  \caption{Superficie de demanda de transporte estimada (viajes por día por \ageb,
    escala logarítmica). La demanda se concentra en las zonas densas de menor
    propiedad vehicular del centro y el oriente.}
  \label{fig:demand}
\end{figure}

\section{Calibración del modelo gravitacional (W2)}
\label{sec:res-w2}

Calibrar el parámetro de decaimiento contra la Encuesta Origen--Destino 2022
(\cref{eq:calibration}) selecciona $\beta^\star = \num{1.2005}$, un decaimiento
notablemente más débil que el prior inicial $\beta = \num{2.0}$: los viajeros de \zmg
recorren distancias mayores de lo que suponía el prior. De manera crucial, el valor
calibrado \emph{mejora} el ajuste en el universo corregido de \num{1881} \ageb{}s en
todas las métricas (\cref{tab:calibration}), revirtiendo la conclusión provisional
obtenida sobre la tabla de \ageb{}s previa, mal filtrada. El $\beta$ calibrado se
adopta para la superficie de demanda y todas las etapas posteriores.

\begin{table}[htbp]
\centering
\caption{Calibración del modelo gravitacional contra la encuesta EOD 2022.}
\label{tab:calibration}
\begin{tabular}{lS[table-format=4.1]S[table-format=1.4]}
\toprule
Modelo & {RMSE (viajes)} & {$R^2$} \\
\midrule
Prior, $\beta = 2.0$        & 5088.2 & {--} \\
Calibrado, $\beta = 1.2005$ & 4524.9 & 0.2498 \\
\bottomrule
\end{tabular}
\end{table}

\section{Diagnóstico de brecha de cobertura (W3)}
\label{sec:res-w3}

La accesibilidad de oportunidades acumuladas (\cref{eq:accessibility}) combinada con
la superficie de demanda arroja el índice de brecha de cobertura (\cref{eq:gap}).
Clasificar los \ageb{}s por quintiles conjuntos de demanda/accesibilidad da
\num{390} de brecha alta (\SI{20.7}{\percent}), \num{1413} de brecha media y \num{78}
de brecha baja (\cref{tab:gap}). Como mostró la \cref{fig:coverage-gap}, las zonas de
brecha alta forman una periferia alrededor de un núcleo bien servido (brecha baja).

\begin{table}[htbp]
\centering
\caption{Clasificación de brecha de cobertura de los \num{1881} \ageb{}s de \zmg.}
\label{tab:gap}
\begin{tabular}{lS[table-format=4.0]S[table-format=2.1]}
\toprule
Categoría & {\ageb{}s} & {Proporción (\%)} \\
\midrule
Brecha alta  & 390  & 20.7 \\
Brecha media & 1413 & 75.1 \\
Brecha baja  & 78   & 4.1 \\
\bottomrule
\end{tabular}
\end{table}

Entrenar un modelo supervisado sobre el objetivo binario de brecha alta
(\cref{sec:w3}), con todas las variables de oferta excluidas, recupera bien la
estructura de la brecha y sin fuga de información (\cref{tab:retrain}). La atribución
SHAP ordena la población, el rezago social (\code{pe\_rezago}) y el proxy de empleo
como los impulsores dominantes: las zonas de brecha alta son áreas densas, de alta
necesidad y ricas en empleo que la red actual atiende de forma insuficiente.

\begin{table}[htbp]
\centering
\caption{Reentrenamiento supervisado sobre la brecha de cobertura (partición de prueba).}
\label{tab:retrain}
\begin{tabular}{lS[table-format=1.3]S[table-format=1.3]}
\toprule
Modelo & {PR-AUC} & {ROC-AUC} \\
\midrule
Random Forest & 0.877 & 0.962 \\
LightGBM      & 0.883 & 0.962 \\
\bottomrule
\end{tabular}
\end{table}

\section{Priorización Nodo--Lugar--Personas (W4)}
\label{sec:res-w4}

El procedimiento CRITIC$\,\times\,$EWM (\cref{eq:critic}) asigna los mayores pesos
objetivos al proxy de empleo (\num{0.163}), la densidad de servicios (\num{0.145}) y
la densidad de intersecciones de cuatro ramales (\num{0.097}); el vector completo de
pesos aparece en la \cref{fig:weights}. Promediada sobre todos los \ageb{}s, la
puntuación de lugar es $\overline{\mathrm{NPP}} = \num{0.4595}$, el término de
equidad $\overline{\mathrm{Eq}} = \num{0.2274}$ y el compuesto
$\overline{F} = \num{0.4131}$ con $\alpha = \num{0.20}$; la media baja de equidad es
esperada, ya que \zmg es predominantemente de baja marginación. Un análisis de
sensibilidad sobre $\alpha \in \{0.10, 0.20, 0.30\}$ confirma que la priorización es
robusta: el ordenamiento correlaciona con la línea base de $\alpha = \num{0.20}$ en
Spearman $\ge \num{0.98}$, de modo que el peso de equidad cambia cuáles \ageb{}s
específicos encabezan la lista sin alterar la estructura general.

\section{Corredores generados (W6)}
\label{sec:res-w6}

El generador produce cinco corredores candidatos, de los cuales cuatro son factibles
(\cref{tab:corridors}); el quinto (\code{W6\_G05}) se rechaza por la restricción de
rectitud por anclas (\num{1.93} > \num{1.8}). La \cref{fig:corridors} mapea los cuatro
corredores factibles sobre la superficie de brecha, y la \cref{fig:pareto} sitúa a
todos los candidatos en el espacio de objetivos.

\begin{table}[htbp]
\centering
\caption{Corredores candidatos generados. \code{W6\_G05} incumple la restricción de
  rectitud por anclas.}
\label{tab:corridors}
\begin{tabular}{l S[table-format=2.1] S[table-format=2.0] S[table-format=6.0] l c}
\toprule
Corredor & {km} & {\ageb{}s} & {Demanda/día} & Modo & Factible \\
\midrule
W6\_G00 & 7.3  & 18 & 66041  & BRT         & Sí \\
W6\_G01 & 23.0 & 27 & 96839  & Tren ligero & Sí \\
W6\_G02 & 12.1 & 25 & 192357 & Tren ligero & Sí \\
W6\_G03 & 2.4  & 5  & 35784  & BRT         & Sí \\
W6\_G05 & 5.4  & 14 & 80234  & Tren ligero & {--} \\
\bottomrule
\end{tabular}
\end{table}

\begin{figure}[htbp]
  \centering
  \includegraphics[width=0.78\linewidth]{w6_pareto.pdf}
  \caption{Corredores candidatos en el espacio de objetivos $(f_2, f_1)$, coloreados
    por factibilidad. Los corredores factibles cubren un rango de compromisos
    longitud/ganancia.}
  \label{fig:pareto}
\end{figure}

El resultado más destacado es \code{W6\_G02}: un corredor sustantivo de
\SI{12.1}{\kilo\metre} y \num{25} \ageb{}s que transporta una estimación de
\num{192357} viajes/día. En las pruebas de mérito introducidas para la validación
(\cref{sec:res-w8}) es el único corredor factible que pasa las tres: el
\SI{56}{\percent} de sus \ageb{}s servidos son de brecha alta (frente a la línea base
metropolitana de \SI{20.7}{\percent}), no es redundante con las rutas existentes
(mejor Jaccard \num{0.18}) y su demanda por kilómetro se ubica en el percentil
\num{73} de la red existente. \code{W6\_G03} también pasa, pero es un tramo corto de
\SI{2.4}{\kilo\metre}, mientras que \code{W6\_G00} y \code{W6\_G01} atienden necesidad
real pero son conectores de menor eficiencia ---un resultado mixto pero caracterizado
honestamente que se discute en el \cref{ch:discusion}.

\section{Auditoría de rutas existentes (W7)}
\label{sec:res-w7}

Evaluar las \num{247} rutas existentes de SITEUR con la misma función multiobjetivo
marca \num{229} de ellas (\cref{tab:audit}): \num{109} como indirectas, \num{80} como
de baja demanda y \num{40} como redundantes, dejando solo \num{18} sin marcar. La
\cref{fig:w7audit} grafica cada ruta en el espacio de objetivos por marca. Cada ruta
marcada lleva una propuesta de modificación (atajo, fusión o retiro), lo que hace la
auditoría directamente comparable con los corredores generados.

\begin{table}[htbp]
\centering
\caption{Marcas de auditoría en las \num{247} rutas existentes de SITEUR.}
\label{tab:audit}
\begin{tabular}{lS[table-format=3.0]}
\toprule
Marca & {Rutas} \\
\midrule
Indirecta    & 109 \\
Baja demanda & 80 \\
Redundante   & 40 \\
Sin marcar   & 18 \\
\bottomrule
\end{tabular}
\end{table}

\begin{figure}[htbp]
  \centering
  \includegraphics[width=0.82\linewidth]{w7_audit.pdf}
  \caption{Auditoría de rutas existentes: las \num{247} rutas de SITEUR en el espacio
    $(f_2, f_1)$, coloreadas por marca de auditoría.}
  \label{fig:w7audit}
\end{figure}

\section{Validación (W8)}
\label{sec:res-w8}

\paragraph{Retro-prueba de exclusión y comparación.} Enmascarar el nivel premium (Mi
Macro y Mi Tren) y volver a ejecutar el generador reproduce las rutas enmascaradas con
un traslape medio de trazado de apenas \SI{15.0}{\percent}; enmascarar líneas férreas
individuales da \SI{0}{\percent}. La comparación contra \num{33} rutas premium
existentes da un traslape medio de \SI{10.5}{\percent} para \SI{44.9}{\kilo\metre} de
corredor generado. El traslape bajo es el resultado esperado y deseado: el generador
identifica zonas genuinamente desatendidas en lugar de replicar líneas existentes. La
reconstrucción es, en todo caso, un proxy débil del mérito de un corredor
(\cref{ch:discusion}).

\paragraph{Métricas antes/después.} Agregar los corredores factibles a la red
existente mejora todas las métricas de cobertura (\cref{tab:beforeafter}): la tasa de
cobertura de \ageb{}s sube \num{1.1} puntos porcentuales, el Gini de accesibilidad cae
\num{0.0187} y \num{47} \ageb{}s antes desatendidos ---hogar de \num{120648}
personas--- ganan servicio.

\begin{table}[htbp]
\centering
\caption{Métricas de cobertura antes/después para los corredores W6 factibles.}
\label{tab:beforeafter}
\begin{tabular}{lS[table-format=2.1]S[table-format=2.1]S[table-format=+1.4]}
\toprule
Métrica & {Antes} & {Después} & {$\Delta$} \\
\midrule
Tasa de cobertura (\%) & 69.9   & 71.0   & +1.1 \\
Gini de accesibilidad  & 0.6333 & 0.6146 & -0.0187 \\
\bottomrule
\end{tabular}
\end{table}

\paragraph{Corroboración fuera de muestra.} La validación más fuerte es externa: la
instantánea \textsc{gtfs} de 2024 es anterior a la apertura en diciembre de 2025 de la
Línea~4 de SITEUR, por lo que el diagnóstico trata el corredor de la Línea~4 como
desatendido. Sus \ageb{}s muestran \num{1.6} veces la tasa metropolitana de brecha
alta, confirmando fuera de muestra que el índice de brecha de cobertura señala
corredores que los planificadores decidieron construir de forma independiente. Esto
corrobora la capa de \emph{diagnóstico}; no prueba que el generador hubiera producido
la Línea~4, distinción que se desarrolla en el \cref{ch:discusion}.
